% Chapter 17: The Lawson Conjecture
\let\LocalMacro\relax

\chapter{Lawson Conjecture}

\section{The Problem}

\subsection{Informal}
A soap film stretched across a wire frame naturally pulls itself as tight as possible at every point, curving equally in opposite directions like a saddle. If you shape the frame like a donut, the soap film forms a beautifully symmetric shape. The Lawson conjecture, proposed in 1970, asks: is this standard soap-film donut the only one that never crosses through itself, or are there other, more exotic possibilities hiding out there?

\subsection{Formal}
\begin{conjecture}[Lawson, 1970]
The only embedded minimal torus in $S^3$ is the Clifford torus (up to rigid motions).
\end{conjecture}

\subsection{Historical Context}
Lawson (1970) classified all minimal tori in $S^3$ as a one-parameter family. The Clifford torus is the unique embedded member, but proving that no other embedded minimal torus exists resisted all attempts for over 40 years.

\subsection{What Was Known}
Lawson constructed a family of minimal surfaces in 1970 and conjectured that his torus was the only embedded minimal torus of genus 1 in the 3-sphere. Partial results were obtained over the following decades, but the general case was open for over 40 years. Brendle's 2013 proof used innovative techniques from geometric measure theory.

\subsection{Why It Matters}
Soap films are nature's solution to a deep optimization problem: they minimize area while spanning a given boundary. The Lawson conjecture asks whether the most symmetric soap-film donut is the only one of its kind. Understanding this question sheds light on the classification of minimal surfaces, which are fundamental objects in differential geometry and mathematical physics.

\subsection{Simplified Version}
Lawson himself classified all minimal tori in $S^3$ as a one-parameter family in 1970. Within this family, the Clifford torus is the unique \emph{embedded} (self-avoiding) member---all others self-intersect. The conjecture asks whether this is true more generally: is the Clifford torus the only embedded minimal torus in $S^3$, regardless of whether it comes from Lawson's family?

\section{Mathematical Theory}


\subsection{Prerequisites}
This chapter assumes familiarity with:
\begin{itemize}
\item Differential geometry of surfaces (minimal surfaces, second fundamental form)
\item Riemannian geometry (curvature, Jacobi fields)
\item Maximum principle methods in geometric analysis
\end{itemize}

\subsection{Lawson's Classification}
Lawson showed that minimal tori in $S^3$ form a one-parameter family $\{\Sigma_\theta\}_{\theta \in [0, \pi/2]}$, where:
\begin{itemize}
\item $\Sigma_0 = T_{\text{Cl}}$ is the Clifford torus.
\item For $\theta > 0$, the surfaces $\Sigma_\theta$ are immersed but not embedded (they self-intersect).
\end{itemize}

\subsection{The Key Difficulty}
To prove the Lawson conjecture, one must show that among all minimal tori in $S^3$, only the Clifford torus is embedded. The challenge is that the space of all minimal tori is infinite-dimensional, and one needs a global argument that rules out every possible competitor.

\section{The Solution}

\subsection{The Big Idea}

Lawson classified all minimal tori in $S^3$ as a one-parameter family, and within that family the Clifford torus is obviously the only embedded member. But the conjecture asks something stronger: \emph{could there be a rogue minimal torus, outside Lawson's family, that is also embedded?} Think of it like this: you have a catalog of all known minimal tori, and the Clifford torus stands out as the only non-self-intersecting one. But how do you know you have not missed one?

Brendle's approach was to use Simon's integral identity as a \emph{budget equation}. The identity says the total curvature of the surface equals a topological cost plus a non-negative ``waste'' term. For a torus, the topological cost is zero---the topology is ``free.'' So the entire curvature budget is accounted for by waste. If the curvature exceeds 2 anywhere, the waste must be positive---but Brendle chose a test function so perfectly that when curvature is too high, the waste \emph{cannot} be large enough to balance the budget. Contradiction. The art was not inventing a new tool but finding the right \emph{input} to a well-understood machine.

\subsection{What Was Stopping Progress}

For over 30 years, geometers attacked this from multiple angles and failed. Curvature flow methods could deform surfaces toward minimal ones but could not establish uniqueness among all embedded tori. Variational methods tried to analyze the infinite-dimensional moduli space directly but could not control its global topology. Classification theorems could show uniqueness \emph{within Lawson's family} but could not rule out tori \emph{outside} it. The maximum principle and Simon's integral identity were both well-known tools, and both had been tried before. The subtle problem was that Simon's identity relates curvature to the gradient of a test function $\phi$, and a poor choice of $\phi$ yields a gradient term too large to control. Previous researchers chose test functions that were either too simple or too general---none exploited the specific geometry of tori in $S^3$.

\subsection{The Breakthrough}

Brendle chose \emph{exactly the right test function} for the maximum principle---one tailored to exploit the topology of tori ($\chi = 0$, making the topological term vanish) and the specific geometry of $S^3$. With this choice, if $|A|^2 > 2$ anywhere, the maximum principle forces a strict inequality that contradicts Simon's integral identity. Therefore $|A|^2 = 2$ everywhere, which characterizes the Clifford torus. The proof is remarkably short (about 15 pages) and uses only classical tools.

\subsection{How It Works}

Simon Brendle's 2013 breakthrough was deceptively simple in its final form: he chose \emph{exactly the right test function} for the maximum principle. The novelty was:

\begin{enumerate}
\item \textbf{A tailored test function}: Brendle constructed a function $\phi$ built from the second fundamental form $|A|^2$ and the geometry of the torus, specifically designed so that the maximum principle applied to $\phi$ (not to $|A|^2$ directly) yields a contradiction if $|A|^2 > 2$ at any point.

\item \textbf{Exploiting the torus topology}: The key insight was that tori have $\chi(\Sigma) = 0$, so the topological term in Simon's identity vanishes. This means the integral identity is \emph{exactly} an identity about $\int |A|^2\,dA$ and the gradient of $\phi$, with no leftover topological correction. Brendle's test function was chosen to make the gradient term $\int |\nabla\phi|^2\,dA$ precisely measurable against the curvature term.

\item \textbf{Rigidity at equality}: If $|A|^2 > 2$ somewhere, the maximum principle forces a strict inequality in Simon's identity, contradicting the integral identity. Therefore $|A|^2 = 2$ everywhere, which is the defining property of the Clifford torus.
\end{enumerate}

Think of Simon's identity as a \emph{budget equation}:

\[
\underbrace{\int_\Sigma (|A|^2 - 2)\,dA}_{\text{``curvature budget''}} = \underbrace{4\pi \chi(\Sigma)}_{\text{topological cost}} + \underbrace{\int_\Sigma |\nabla \phi|^2\,dA}_{\text{``waste term''} \ge 0}
\]

For a torus, the topological cost is zero (since $\chi = 0$). So the equation says: \emph{the total curvature budget equals the waste}, and the waste is always non-negative. If $|A|^2 > 2$ anywhere, the left side is positive, so the waste must be positive---but Brendle's choice of $\phi$ shows that if $|A|^2$ achieves a large enough maximum, the waste term \emph{cannot} be large enough to balance. This is a contradiction.

\begin{theorem}[Brendle \cite{brendle2013}]
The only embedded minimal torus in $S^3$ is the Clifford torus (up to rigid motions).
\end{theorem}

Brendle's proof proceeds in three steps:

\begin{enumerate}
\item \textbf{Simon's identity}: A sharp integral inequality relating the second fundamental form $|A|^2$ to the metric:
\[
\int_\Sigma \left( |A|^2 - 2 \right) \, dA = 4\pi \chi(\Sigma) + \int_\Sigma |\nabla \phi|^2 \, dA
\]
where $\phi$ is a specific function constructed from the second fundamental form.

\item \textbf{Maximum principle}: A clever application of the Omori-Yau maximum principle to the quantity $|A|^2$, showing that if $|A|^2 > 2$ at any point, then $|A|^2$ must achieve a maximum that contradicts Simon's identity.

\item \textbf{Rigidity}: Equality in Simon's identity forces $|A|^2 = 2$ everywhere, which characterizes the Clifford torus.
\end{enumerate}

\begin{highlightbox}
\textbf{The punchline:} Brendle's proof is remarkably short (about 15 pages) and uses only classical differential geometry---Simon's integral identity, the maximum principle, and rigidity. The breakthrough was not a new theorem or a new technique, but the \emph{correct instantiation} of an existing framework: the right test function, exploiting the topology of tori, to close a 30-year gap. This is a powerful lesson: sometimes the hardest part of a proof is choosing the right input to a well-understood machine.
\end{highlightbox}

\subsection{Brendle's Proof in Detail}

\begin{theorem}[Brendle \cite{brendle2013}]
The only embedded minimal torus in $S^3$ is the Clifford torus (up to rigid motions).
\end{theorem}

Brendle's proof proceeds in three steps:

\begin{enumerate}
\item \textbf{Simon's identity}: A sharp integral inequality relating the second fundamental form $|A|^2$ to the metric:
\[
\int_\Sigma \left( |A|^2 - 2 \right) \, dA = 4\pi \chi(\Sigma) + \int_\Sigma |\nabla \phi|^2 \, dA
\]
where $\phi$ is a specific function constructed from the second fundamental form.

\item \textbf{Maximum principle}: A clever application of the Omori-Yau maximum principle to the quantity $|A|^2$, showing that if $|A|^2 > 2$ at any point, then $|A|^2$ must achieve a maximum that contradicts Simon's identity.

\item \textbf{Rigidity}: Equality in Simon's identity forces $|A|^2 = 2$ everywhere, which characterizes the Clifford torus.
\end{enumerate}

\subsection{After the Proof}

The Lawson conjecture is completely settled: Brendle proved that the Clifford torus is the unique embedded minimal torus in $S^3$. The proof uses only classical tools---Simon's integral identity and the maximum principle---and leaves no remaining cases. The specificity of the result (only one ambient space, only one genus) reflects the method rather than any gap in the theorem.

The proof methods, while classical, were applied in a novel way. Simon's identity provides a sharp integral inequality, and the maximum principle converts pointwise information into global conclusions. The specific choice of test function---tailored to exploit the topology of tori and the geometry of $S^3$---was the key innovation. This combination of integral identities with maximum principle arguments has been influential in geometric analysis, demonstrating that the right test function can unlock problems that have resisted more sophisticated machinery.

Several natural generalizations remain open. Classifying minimal surfaces of higher genus in $S^3$ is a major unsolved problem---for genus 2 and above, very little is known. The Lawson conjecture for other ambient spaces, such as other three-dimensional space forms (hyperbolic space, lens spaces), is also largely open. Whether the Clifford torus is the unique embedded minimal torus in other positively curved three-manifolds is an active area of research. Extending Brendle's technique to these settings, or finding fundamentally new approaches, is an ongoing challenge in minimal surface theory.

\section{Impact}
\begin{itemize}
\item \textbf{Minimal surface theory}: Established the Clifford torus as the unique minimal torus in $S^3$, completing Lawson's 1970 classification.
\item \textbf{Maximum principle methods}: Demonstrated the power of combining integral identities with pointwise maximum arguments.
\item \textbf{Connection to Willmore}: The Lawson result complements the Willmore conjecture (Chapter 16) by characterizing the unique minimizer.
\end{itemize}

\section{Exercises}

\begin{exercise}[Basic]
Verify that the Clifford torus is minimal in $S^3$ by computing its mean curvature.
\end{exercise}

\begin{exercise}[Intermediate]
Explain why Lawson's one-parameter family $\Sigma_\theta$ for $\theta > 0$ consists of immersed (not embedded) tori. What causes the self-intersection?
\end{exercise}

\begin{exercise}[Challenge]
Sketch how Simon's identity relates the total curvature $\int |A|^2 \, dA$ to the topology $\chi(\Sigma)$ of the surface. Why is this useful for rigidity?
\end{exercise}

\section{Timeline}

\begin{tabularx}{\linewidth}{lX}
\toprule
\textbf{Year} & \textbf{Event} \\ \midrule
1970 & Lawson classifies minimal tori in $S^3$ \\ 
1993 & Simon's identity proves existence of surfaces minimizing Willmore-type functionals \\
2013 & Brendle proves the Lawson conjecture \\
2014 & Brendle wins Fields Medal \\ \bottomrule
\end{tabularx}

\section{Citations}

\begin{enumerate}
\item Brendle, S. (2013). ``Embedded minimal tori in $S^3$ and the Lawson conjecture.'' Acta Mathematica, 211(1), 177--190.
\item Lawson, H. B. (1970). ``Complete minimal surfaces in $S^3$.'' Annals of Mathematics, 92(2), 335--374.
\item Simon, L. (1993). ``Existence of surfaces minimizing the Willmore functional.'' Communications in Analysis and Geometry, 1(2), 281--326.
\end{enumerate}
